\documentclass[11pt,a4paper]{article}
\usepackage[utf8]{inputenc}
\usepackage{geometry}
\geometry{a4paper, margin=1in}
\usepackage{graphicx}
\usepackage{hyperref}
\usepackage{listings}
\usepackage{xcolor}
\usepackage{amsmath}
\usepackage{tabularx}
\usepackage{booktabs}
\usepackage{float}
\usepackage{tcolorbox}
\usepackage{parskip}

\definecolor{codegreen}{rgb}{0,0.6,0}
\definecolor{codegray}{rgb}{0.5,0.5,0.5}
\definecolor{codepurple}{rgb}{0.58,0,0.82}
\definecolor{backcolour}{rgb}{0.95,0.95,0.92}

\lstdefinestyle{mystyle}{
    backgroundcolor=\color{backcolour},   
    commentstyle=\color{codegreen},
    keywordstyle=\color{blue},
    numberstyle=\tiny\color{codegray},
    stringstyle=\color{codepurple},
    basicstyle=\ttfamily\scriptsize,
    breakatwhitespace=false,         
    breaklines=true,                 
    captionpos=b,                    
    keepspaces=true,                 
    numbers=left,                    
    numbersep=5pt,                  
    showspaces=false,                
    showstringspaces=false,
    showtabs=false,                  
    tabsize=2
}
\lstset{style=mystyle}

\title{\Huge \textbf{Ultimate Code \& System Architecture Guide} \\ \vspace{0.5cm} \Large Smart Air Guardian}
\author{Complete Technical Deep-Dive}
\date{\today}

\begin{document}

\maketitle
\tableofcontents
\newpage

\section{Introduction}
This document serves as the absolute, exhaustive, line-by-line technical manual for the Smart Air Guardian. It covers exactly how the ESP32 code works, how the Machine Learning is translated from Python to C++, and how the PWA Web Dashboard syncs with the system.

\section{Part 1: The Machine Learning Pipeline (Python to C++)}
The intelligence of the system begins offline. We cannot run Python or heavy AI frameworks on an ESP32. Instead, we train models on a computer and export the mathematical constants.

\subsection{Data Generation and Targets}
The script \texttt{train\_final\_models.py} uses historical AQI and weather data. It drops missing values and sets up three targets for the AI to learn:
\begin{itemize}
    \item \textbf{Target 1 (Regression):} Predict the AQI 1 hour into the future.
    \item \textbf{Target 2 (Logistic):} Predict if AQI will exceed 150 (Binary: 0 or 1).
    \item \textbf{Target 3 (Decision Tree):} Classify the danger level (0 to 3).
\end{itemize}

\subsection{Min-Max Scaling (Crucial for ESP32)}
Before the AI can process numbers, they must be scaled down to a value between $0.0$ and $1.0$. The Python script calculates the minimum and maximum bounds for PM2.5, CO, Temperature, and Humidity. 
\textbf{Inside the ESP32 Code}, this is replicated using the \texttt{constrain} function:
\begin{lstlisting}[language=C++]
float s_pm25  = constrain((raw_pm25 - 0.40f) / (302.70f - 0.40f), 0.0f, 1.0f);
float s_co    = constrain((raw_co - 104.0f) / (3650.0f - 104.0f), 0.0f, 1.0f);
\end{lstlisting}

\subsection{C++ Linear Regression (Future AQI)}
Once scaled, the ESP32 multiplies the inputs by the weights generated by Python:
\begin{lstlisting}[language=C++]
float predictFutureAQI(float raw_pm25, float raw_co, float raw_temp, float raw_humid) {
    // Scaling done here...
    float future_aqi = (284.2261f * s_pm25) + (82.0077f * s_co) + 
                       (0.0337f * s_temp) + (-0.0661f * s_humid) + (59.7454f);
    return max(0.0f, future_aqi);
}
\end{lstlisting}

\subsection{C++ Logistic Regression (Hazard Probability)}
Logistic regression uses a "Sigmoid" curve to output a percentage probability. The ESP32 calculates $Z$, then applies the curve.
\begin{lstlisting}[language=C++]
float z = (23.5640f * s_pm25) + (-0.3901f * s_co) + (-0.0083f * s_temp) + (-0.0110f * s_humid) - 4.4767f;
float prob = 1.0f / (1.0f + exp(-z)); // The Sigmoid function
\end{lstlisting}

\section{Part 2: ESP32 Firmware Deep-Dive (\texttt{ESP32\_Air\_Quality.ino})}

\subsection{Initialization Phase (\texttt{setup()})}
When the ESP32 turns on, it assigns pins to inputs (sensors) and outputs (LEDs).
\begin{itemize}
    \item \textbf{Pins 34, 35, 32:} Configured as Analog Inputs to read voltages from the MQ135, MQ7, and GP2Y1010 dust sensors.
    \item \textbf{Pins 25, 27, 12:} Configured as Digital Outputs to drive the RGB LED.
\end{itemize}
It then initializes the I2C bus (\texttt{Wire.begin(21, 22)}) to communicate with the 1.3" OLED display. Finally, it starts a local Web Server on port 80.

\subsection{The Main Loop (\texttt{loop()})}
The loop runs continuously. To prevent locking up, it uses a time-check mechanism (\texttt{millis()}) to read sensors exactly every 1.5 seconds.

\subsubsection{Hardware Oversampling \& Exponential Smoothing}
To ensure the dashboard doesn't display jittery numbers, the code uses aggressive smoothing.
\begin{lstlisting}[language=C++]
// Oversampling: Read gas sensors 24 times rapidly and average them
long sumMQ135 = 0;
for (int i = 0; i < 24; i++) {
    sumMQ135 += analogRead(MQ135_PIN);
    delayMicroseconds(250);
}
int liveRawMQ135 = sumMQ135 / 24;

// Exponential Smoothing: Blend old data (85%) with new data (15%)
liveCO2_ppm = (liveCO2_ppm * 0.85f) + (calculatedCO2 * 0.15f);
\end{lstlisting}

\subsubsection{EPA AQI Calculation}
The code takes the smoothed PM2.5 value and runs it through the official US EPA piecewise linear function to convert micrograms into the 0-500 standard AQI scale.

\subsubsection{Actuator Control (Lights and Sound)}
Based on the Decision Tree's output class, \texttt{setRGBStatus()} switches the LED to Green, Yellow, Red, or Purple. If the Hazard Probability is $\geq 0.85$, it triggers the buzzer.

\section{Part 3: Cloud Synchronization \& Firebase}
The ESP32 does not use a heavy Firebase library. Instead, it natively sends raw HTTPS REST commands, which saves massive amounts of memory.

\subsection{Telemetry Push (Every 3 seconds)}
The ESP32 constructs a JSON string containing the AQI, Temp, Humidity, and Predictions. It uses \texttt{HTTPClient} to send an HTTP \texttt{PUT} request to \texttt{https://air-guardian.../telemetry.json}.

\subsection{Command Polling (Every 1.5 seconds)}
The ESP32 sends a \texttt{GET} request to \texttt{/commands/lastCommand.json}.
If the user clicks "Turn on Buzzer" on the web dashboard, Firebase updates this JSON. The ESP32 parses the string, checks the timestamp (to ensure it doesn't run the same command twice), and executes the physical action.

\section{Part 4: The Web Dashboard (PWA)}
The file \texttt{frontend\_dashboard.html} is the graphical user interface.
\begin{itemize}
    \item \textbf{CSS Variables:} It uses modern CSS variables (e.g., \texttt{--bg-base: \#060b18}) to enforce a strict "Claymorphic" design scheme.
    \item \textbf{JavaScript Logic:} It fetches data from Firebase in real-time. When data arrives, it updates the HTML DOM elements (like the big AQI gauge) instantly.
    \item \textbf{C2 Controls:} When you click a switch on the website, Javascript sends a \texttt{PUT} request to Firebase, which the ESP32 will pick up on its next polling cycle.
\end{itemize}

\end{document}
